\Article{233}
Wenn eine (primitive) Form $F=(a,b,c)$ so beschaffen ist, dass sich zwei Zahlen $g$, $h$ mit $gg\equiv a$, $gh\equiv b$, $hh\equiv c$ nach einem gegebenen Modul $m$ finden lassen, nennen wir diese Form einen quadratischen Rest der Zahl $m$ und $gx+hy$ einen Wert des Ausdrucks $\surd(axx+2bxy+cyy)(\mathrm{mod.}\ m)$, oder kürzer $(g,h)$ einen Wert des Ausdrucks $\surd(a,b,c)$ oder $\surd F(\mathrm{mod.}\ m)$. Allgemeiner: Wenn ein Multiplikator $M$, der zum Modul $m$ teilerfremd ist, so beschaffen ist, dass
\[
gg\equiv aM,\quad gh\equiv bM,\quad hh\equiv cM\quad (\mathrm{mod.}\ m)
\]
gelten kann, nennen wir $M\times(a,b,c)$ oder $MF$ einen quadratischen Rest von $m$ und $(g,h)$ einen Wert des Ausdrucks $\surd M(a,b,c)$ oder $\surd MF(\mathrm{mod.}\ m)$. So ist \textit{zum Beispiel} die Form $(3,1,54)$ ein quadratischer Rest von $23$ und $(7,10)$ ein Wert des Ausdrucks $\surd(3,1,54)(\mathrm{mod.}\ 23)$; ebenso ist $(2,-4)$ ein Wert des Ausdrucks $\surd5(10,3,17)(\mathrm{mod.}\ 23)$. Wozu diese Definitionen dienen, wird sich später zeigen; hier halten wir die folgenden Sätze fest:

\par
\SourceLeaf{242}
I. Ist $M(a,b,c)$ ein quadratischer Rest der Zahl $m$, so teilt diese Zahl die Determinante der Form $(a,b,c)$. Denn ist $(g,h)$ ein Wert des Ausdrucks $\surd M(a,b,c)$ (mod.\ $m$), gilt also
\[
gg\equiv aM,\quad gh\equiv bM,\quad hh\equiv cM\ (\mathrm{mod.}\ m)
\]
so ist $bbMM-acMM\equiv0$, oder $(bb-ac)MM$ ist durch $m$ teilbar. Da aber $M$ als teilerfremd zu $m$ vorausgesetzt wird, ist auch $bb-ac$ durch $m$ teilbar.

II. Ist $M(a,b,c)$ ein quadratischer Rest von $m$ und $m$ entweder eine Primzahl oder eine Primzahlpotenz, etwa $=p^\mu$, so ist der Einzelcharakter der Form $(a,b,c)$ bezüglich der Zahl $p$ entweder $Rp$ oder $Np$, je nachdem, ob $M$ Rest oder Nichtrest von $p$ ist. Das folgt sofort daraus, dass sowohl $aM$ als auch $cM$ Rest von $m$ oder von $p$ ist und mindestens eine der Zahlen $a$, $c$ nicht durch $p$ teilbar ist (Artikel 230).

Ebenso ist, wenn (unter sonst gleichen Bedingungen) $m=4$ gilt, entweder $1,4$ oder $3,4$ der Einzelcharakter der Form $(a,b,c)$, je nachdem, ob $M\equiv1$ oder $\equiv3$ ist; Entsprechendes gilt, wenn $m=8$ oder eine höhere Potenz der Zahl $2$ ist. Dann ist $1,8;\ 3,8;\ 5,8;\ 7,8$ der Einzelcharakter der Form $(a,b,c)$, je nachdem, ob $M\equiv1;\ 3;\ 5;\ 7\ (\mathrm{mod.}\ 8)$ gilt.

III. Umgekehrt: Ist $m$ eine Primzahl oder eine Potenz einer ungeraden Primzahl $=p^\mu$, die die Determinante $bb-ac$ teilt, und ist $M$ Rest oder Nichtrest von $p$, je nachdem, ob der Charakter der Form $(a,b,c)$ bezüglich $p$ gleich $Rp$ oder $Np$ ist, so ist $M(a,b,c)$ ein quadratischer Rest von $m$. Denn wenn $a$ nicht durch $p$ teilbar ist, ist $aM$ Rest von $p$ und damit auch von $m$; ist also $g$ ein Wert des Ausdrucks $\surd aM$ (mod.\ $m$) und $h$ ein Wert des Ausdrucks $\frac{bg}{a}$ (mod.\ $m$), so ist $gg\equiv aM;\ ah\equiv bg$ und damit
\[
agh\equiv bgg\equiv abM\quad\text{und}\quad gh\equiv bM
\]
und schließlich
\[
ahh\equiv bgh\equiv bbM\equiv bbM-(bb-ac)M\equiv acM
\]
und damit $hh\equiv cM$, \textit{d.h.}, $(g,h)$ ist ein Wert des Ausdrucks $\surd M(a,b,c)$. Ist aber $a$ durch $m$ teilbar, so ist $c$ sicher nicht dadurch teilbar; daraus sieht man leicht, dass dasselbe folgt, wenn man für $h$ einen Wert des Ausdrucks $\surd cM$ (mod.\ $m$) nimmt. Für $g$ nimmt man einen Wert des Ausdrucks $\frac{bh}{c}$ (mod.\ $m$).

Ebenso beweist man: Gilt $m$ $=4$ und teilt diese Zahl $bb-ac$, und wählt man die Zahl $M$ entweder $\equiv1$ oder $\equiv3$, je nachdem, ob $1,4$ oder $3,4$ der Einzelcharakter
\SourceLeaf{243}
der Form $(a,b,c)$ ist, so ist $M(a,b,c)$ ein quadratischer Rest von $m$. Ebenso gilt: Gilt $m$ $=8$ oder ist es eine höhere Potenz von $2$, durch die $bb-ac$ teilbar ist, und wählt man $M$ mit $\equiv1;\ 3;\ 5;\ 7\ (\mathrm{mod.}\ 8)$, wie es der Einzelcharakter der Form $(a,b,c)$ bezüglich der Zahl $8$ verlangt, so ist $M(a,b,c)$ ein quadratischer Rest von $m$.

IV. Ist die Determinante der Form $(a,b,c)$ durch $=D$ gegeben und $M(a,b,c)$ ein quadratischer Rest von $D$, so lassen sich alle Einzelcharaktere der Form $(a,b,c)$, sowohl bezüglich der einzelnen ungeraden Primteiler von $D$ als auch bezüglich der Zahl $4$ oder der Zahl $8$ (wenn sie $D$ teilen), sofort aus der Zahl $M$ ablesen. So sind \textit{zum Beispiel}, weil $3(20,10,27)$ ein quadratischer Rest von $440$, nämlich $(150,9)$ ein Wert des Ausdrucks $\surd3(20,10,27)$ nach dem Modul $440$ ist und $3N5$, $3R11$ gilt, die Charaktere der Form $(20,10,27)$ gleich $3,8;\ N5;\ R11$. Nur die Einzelcharaktere bezüglich der Zahlen $4$ und $8$ haben, wenn diese die Determinante nicht teilen, keinen notwendigen Zusammenhang mit der Zahl $M$

V. Umgekehrt: Enthält die Zahl $M$, die zu $D$ teilerfremd ist, alle Einzelcharaktere der Form $(a,b,c)$ (mit Ausnahme der Charaktere bezüglich der Zahlen $4,8$, wenn diese $D$ nicht teilen), so ist $M(a,b,c)$ ein quadratischer Rest von $D$. Denn aus III folgt: Bringt man $D$ auf die Form $\pm A^\alpha B^\beta D^\gamma\ldots$, wobei $A$, $B$, $C$ usw. verschiedene Primzahlen sind, so ist $M(a,b,c)$ ein quadratischer Rest jeder der Zahlen $A^\alpha$, $B^\beta$, $C^\gamma$ usw. Ist also ein Wert des Ausdrucks $\surd M(a,b,c)$ nach dem Modul $A^\alpha$ gleich $(\mathfrak{A},\mathfrak{A}')$; nach dem Modul $B^\beta$ gleich $(\mathfrak{B},\mathfrak{B}')$; nach dem Modul $C^\gamma$ gleich $(\mathfrak{C},\mathfrak{C}')$ usw., und bestimmt man die Zahlen $g$, $h$ so, dass $g\equiv\mathfrak{A},\mathfrak{B},\mathfrak{C}$ usw.; $h\equiv\mathfrak{A}',\mathfrak{B}',\mathfrak{C}'$ usw. nach den Moduln $A^\alpha$, $B^\beta$, $C^\gamma$ usw. gilt (Artikel 32), so sieht man leicht, dass $gg\equiv aM$, $gh\equiv bM$, $hh\equiv cM$ nach allen Moduln $A^\alpha$, $B^\beta$, $C^\gamma$ usw. gilt und damit auch nach dem Modul $D$, der ihr Produkt ist.

VI. Aus diesen Gründen nennen wir Zahlen wie $M$ \textit{charakteristische Zahlen} der Form $(a,b,c)$. Nach V lassen sich mehrere solche Zahlen ohne Schwierigkeit finden, sobald alle Einzelcharaktere dieser Form bekannt sind; die einfachsten findet man aber meist am leichtesten durch Probieren. Offenbar sind, wenn $M$ eine charakteristische Zahl einer gegebenen primitiven Form mit der Determinante $D$ ist, alle Zahlen, die $M$ nach dem Modul $D$ kongruent sind, charakteristische Zahlen derselben Form; Formen in derselben Klasse oder auch in verschiedenen Klassen desselben Geschlechts haben dieselben charakteristischen Zahlen, weshalb jede charakteristische
\SourceLeaf{244}
Zahl einer gegebenen Form auch der ganzen Klasse und dem ganzen Geschlecht zugeordnet werden kann; schließlich ist $1$ stets eine charakteristische Zahl der Hauptform, Hauptklasse und des Hauptgeschlechts, oder jede Form des Hauptgeschlechts ist ein Rest ihrer eigenen Determinante.

VII. Ist $(g,h)$ ein Wert des Ausdrucks $\surd M(a,b,c)$ (mod.\ $m$) und gilt $g'\equiv g$, $h'\equiv h$ (mod.\ $m$), so ist auch $(g',h')$ ein Wert desselben Ausdrucks. Solche Werte können als \textit{äquivalent} gelten; sind dagegen $(g,h)$, $(g',h')$ Werte desselben Ausdrucks $\surd M(a,b,c)$, ohne dass zugleich $g'\equiv g$, $h'\equiv h$ (mod.\ $m$) gilt, so sind sie als \textit{verschieden} zu zählen. Offenbar ist, wenn $(g,h)$ ein Wert eines solchen Ausdrucks ist, auch $(-g,-h)$ ein solcher Wert, und man beweist leicht, dass diese Werte stets verschieden sind, außer wenn $m=2$ ist. Ebenso leicht beweist man, dass der Ausdruck $\surd M(a,b,c)$ (mod.\ $m$) nicht mehr verschiedene Werte als zwei solche (entgegengesetzte) haben kann, wenn $m$ entweder eine ungerade Primzahl, eine Potenz einer ungeraden Primzahl oder $=4$ ist; gilt dagegen $m$ $=8$ oder ist es eine höhere Potenz der Zahl $2$, so gibt es genau vier. Daraus folgt mit VI leicht: Ist die Determinante $D$ der Form $(a,b,c)$ durch $=\pm2^\mu A^\alpha B^\beta\ldots$ gegeben, wobei $A$, $B$ usw. verschiedene ungerade Primzahlen bezeichnen, deren Anzahl $=n$ ist, und ist $M$ eine charakteristische Zahl dieser Form, so gibt es genau $2^n$, $2^{n+1}$ oder $2^{n+2}$ verschiedene Werte des Ausdrucks $\surd M(a,b,c)$ (mod.\ $D$), je nachdem, ob $\mu$ $<2$, $=2$ oder $>2$ ist. So gibt es \textit{zum Beispiel} sechzehn Werte des Ausdrucks $\surd7(12,6,-17)$ (mod.\ $240$), nämlich $(\pm18,\mp11)$, $(\pm18,\pm29)$, $(\pm18,\mp91)$, $(\pm18,\pm109)$, $(\pm78,\pm19)$, $(\pm78,\pm59)$, $(\pm78,\mp61)$, $(\pm78,\mp101)$. Einen ausführlicheren Beweis lassen wir der Kürze wegen weg, da er für das Folgende nicht besonders nötig ist.

VIII. Schließlich bemerken wir: Haben zwei äquivalente Formen $(a,b,c)$, $(a',b',c')$ die Determinante $D$ und die charakteristische Zahl $M$, und geht die erste durch die Substitution $\alpha,\beta,\gamma,\delta$ in die zweite über, so ergibt sich aus jedem Wert des Ausdrucks $\surd M(a,b,c)$, etwa $(g,h)$, ein Wert des Ausdrucks $\surd M(a',b',c')$, nämlich $(\alpha g+\gamma h,\beta g+\delta h)$. Den Beweis kann jeder ohne Schwierigkeit finden.

\Needspace{12\baselineskip}
\begin{center}
\textit{Komposition der Formen.}
\end{center}

\Needspace{8\baselineskip}
\Article{234}
Nachdem wir die Einteilung der Formen in Klassen, Geschlechter und Ordnungen besprochen und die allgemeinen Eigenschaften erklärt haben, die sich unmittelbar aus diesen Unterscheidungen ergeben, gehen wir zu einem anderen sehr wichtigen Gegenstand über, den bisher noch niemand
\SourceLeaf{245}
behandelt hat: der \textit{Komposition} der Formen. Am Anfang dieser Untersuchung schieben wir gleich das folgende Lemma ein, damit wir später die fortlaufende Folge der Beweise nicht unterbrechen müssen.

\textsc{Lemma.} \textit{Gegeben seien vier Folgen ganzer Zahlen}
\[
\begin{gathered}a,a',a''\ldots a^n;\qquad b,b',b''\ldots b^n;\\c,c',c''\ldots c^n;\qquad d,d',d''\ldots d^n\end{gathered}
\]
\textit{mit gleich vielen (etwa $n+1$) Gliedern, so beschaffen, dass}
\[
cd'-dc',\quad cd''-dc''\ \text{\textit{etc.}},\quad c'd''-d'c''\ \text{\textit{etc. etc.}}
\]
\textit{jeweils gleich}
\[
=k(ab'-ba'),\quad k(ab''-ba'')\ \text{\textit{etc.}},\quad k(a'b''-b'a'')\ \text{\textit{etc. etc.}}
\]
\textit{sind, oder allgemein}
\[
c^\lambda d^\mu-d^\lambda c^\mu=k(a^\lambda b^\mu-b^\lambda a^\mu)
\]
\textit{gilt; dabei bezeichnet $k$ eine gegebene ganze Zahl; $\lambda$, $\mu$ sind beliebige verschiedene ganze Zahlen zwischen $0$ und $n$ einschließlich, deren größere $\mu$ ist}\HistoricalNote{*)}{Dabei betrachten wir $a$ als $a^0$, $b$ als $b^0$ usw. --- Übrigens gilt dieselbe Gleichung offenbar auch, wenn $\lambda=\mu$ oder $\lambda>\mu$ ist.}\textit{; außerdem haben alle $a^\lambda b^\mu-b^\lambda a^\mu$ keinen gemeinsamen Teiler. Dann lassen sich vier ganze Zahlen $\alpha,\beta,\gamma,\delta$ finden, sodass}
\[
\begin{aligned}
\alpha a+\beta b&=c, &\alpha a'+\beta b'&=c', &\alpha a''+\beta b''&=c''\ \text{\textit{etc.}}\\
\gamma a+\delta b&=d, &\gamma a'+\delta b'&=d', &\gamma a''+\delta b''&=d''\ \text{\textit{etc.}}
\end{aligned}
\]
\textit{oder allgemein}
\[
\alpha a^\nu+\beta b^\nu=c^\nu,\quad\gamma a^\nu+\delta b^\nu=a^\nu
\]
\textit{gilt; dann ist}
\[
\alpha\delta-\beta\gamma=k
\]

Da nach Voraussetzung die Zahlen $ab'-ba'$, $ab''-ba''$ usw., $a'b''-b'a''$ usw. (deren Anzahl $=\frac12(n+1)n$ ist) keinen gemeinsamen Teiler haben, lassen sich ebenso viele andere ganze Zahlen finden, sodass die Summe der Produkte $=1$ ist, wenn man jene jeweils mit diesen multipliziert (Artikel 40). Bezeichnen wir diese Multiplikatoren mit $(0,1)$, $(0,2)$ usw., $(1,2)$ usw., oder allgemein den Multiplikator von $a^\lambda b^\mu-b^\lambda a^\mu$ mit $(\lambda,\mu)$, sodass
\[
\Sigma(\lambda,\mu)(a^\lambda b^\mu-b^\lambda a^\mu)=1
\]
gilt. (Mit dem Buchstaben $\Sigma$ bezeichnen wir die Summe aller Werte des Ausdrucks, der auf ihn
\SourceLeaf{246}
folgt, die entstehen, wenn man $\lambda$, $\mu$ alle verschiedenen Werte zwischen $0$ und $n$ mit $\mu>\lambda$ gibt.) Setzen wir nun
\[
\begin{aligned}
\Sigma(\lambda,\mu)(c^\lambda b^\mu-b^\lambda c^\mu)&=\alpha, &\Sigma(\lambda,\mu)(a^\lambda c^\mu-c^\lambda a^\mu)&=\beta\\
\Sigma(\lambda,\mu)(d^\lambda b^\mu-b^\lambda d^\mu)&=\gamma, &\Sigma(\lambda,\mu)(a^\lambda d^\mu-d^\lambda a^\mu)&=\delta
\end{aligned}
\]
so haben diese $\alpha,\beta,\gamma.\delta$ die verlangten Eigenschaften.

\textit{Beweis.} I. Bezeichnet $\nu$ eine beliebige ganze Zahl zwischen $0$ und $n$, so ist
\[
\begin{aligned}
\alpha a^\nu+\beta b^\nu
&=\Sigma(\lambda,\mu)(c^\lambda b^\mu a^\nu-b^\lambda c^\mu a^\nu+a^\lambda c^\mu b^\nu-c^\lambda a^\mu b^\nu)\\
&=\frac1k\Sigma(\lambda,\mu)(c^\lambda d^\mu c^\nu-d^\lambda c^\mu c^\nu)\\
&=\frac1k c^\nu\Sigma(\lambda,\mu)(c^\lambda d^\mu-d^\lambda c^\mu)\\
&=c^\nu\Sigma(\lambda,\mu)(a^\lambda b^\mu-b^\lambda a^\mu)=c^\nu
\end{aligned}
\]
Eine ähnliche Rechnung ergibt
\[
\gamma a^\nu+\delta b^\nu=d^\nu.\qquad\text{\textit{Q. E. P.}}
\]

II. Da also
\[
c^\lambda=\alpha a^\lambda+\beta b^\lambda,\quad c^\mu=\alpha a^\mu+\beta b^\mu
\]
gilt, ist
\[
c^\lambda b^\mu-b^\lambda c^\mu=\alpha(a^\lambda b^\mu-b^\lambda a^\mu)
\]
und ebenso
\[
\begin{aligned}
a^\lambda c^\mu-c^\lambda a^\mu&=\beta(a^\lambda b^\mu-b^\lambda a^\mu)\\
d^\lambda b^\mu-b^\lambda d^\mu&=\gamma(a^\lambda b^\mu-b^\lambda a^\mu)\\
a^\lambda d^\mu-d^\lambda a^\mu&=\delta(a^\lambda b^\mu-b^\lambda a^\mu)
\end{aligned}
\]
Aus diesen Formeln lassen sich die Werte von $\alpha,\beta,\gamma,\delta$ viel leichter finden, wenn man nur $\lambda$, $\mu$ so wählt, dass $a^\lambda b^\mu-b^\lambda a^\mu$ nicht $=0$ ist. Das ist sicher möglich, weil nach Voraussetzung alle $a^\lambda b^\mu-b^\lambda a^\mu$ keinen gemeinsamen Teiler haben und daher nicht \textit{alle} $=0$ sein können. --- Multipliziert man die erste dieser Gleichungen mit der vierten und die zweite mit der dritten und subtrahiert, so ergibt sich aus denselben Gleichungen
\[
(\alpha\delta-\beta\gamma)(a^\lambda b^\mu-b^\lambda a^\mu)^2
=(a^\lambda b^\mu-b^\lambda a^\mu)(c^\lambda d^\mu-d^\lambda c^\mu)
=k(a^\lambda b^\mu-b^\lambda a^\mu)^2
\]
und damit notwendig
\[
\alpha\delta-\beta\gamma=k.\qquad\text{\textit{Q. E. S.}}
\]
